\begin{description}
\item[(T12.1a)] Total energy of any particle is
  $E = E_{\mathrm{grav}} + mc^2$, so we have
  \[
    E_1 = -\frac{G M_{\mathrm{bh}} m}{R_{\mathrm{SC}}} + mc^2, \qquad
    E_2 = -\frac{G M_{\mathrm{bh}} m}{\kappa R_{\mathrm{SC}}} + mc^2.
  \]
  \hfill(1.0)

  We require $E_1 + E_2 = 0$. Hence,
  \[
    2mc^2 - \frac{G M_{\mathrm{bh}} m}{R_{\mathrm{SC}}}
    \left(1 + \frac{1}{\kappa}\right) = 0
  \]
  \hfill(1.0)

  Then
  \[
    \frac{1}{\kappa} = \frac{2 c^2 R_{\mathrm{SC}}}{G M_{\mathrm{bh}}} - 1 = 3
  \]
  where we have used
  $R_{\mathrm{SC}} = \dfrac{2 G M_{\mathrm{bh}}}{c^2}$. \hfill(1.0)

  This gives
  \[
    \kappa = \frac{1}{3}.
  \]
  \hfill(1.0)

\item[(T12.1b)] Using Wien's law,
  $\lambda_{\mathrm{bb}} = b/T_{\mathrm{bh}}$ where
  $b = 2.898 \times 10^{-3}$\,m\,K. \hfill(1.0)

  Hence,
  \[
    T_{\mathrm{bh}} = \frac{b}{16 R_{\mathrm{SC}}}
    \implies T_{\mathrm{bh}} = \frac{b c^2}{32 G M_{\mathrm{bh}}}
  \]
  \hfill(1.0)

  Since $R_{\mathrm{SC}}$ for a solar mass black hole is 2.952\,km, for a
  black hole with mass $= 10\,M_\odot$, we get
  $R_{\mathrm{SC,\,10\odot}} \simeq 29.52$\,km. \hfill(1.0)

  The temperature of this BH will be
  \[
    T_{\mathrm{bh,\,10\odot}}
    = \frac{2.898 \times 10^{-3}\,\mathrm{m\,K}}{16 \times
      R_{\mathrm{SC,\,10\odot}}}
  \]
  \[
    T_{\mathrm{bh,\,10\odot}} = 6.1 \times 10^{-9}\,\mathrm{K}.
  \]
  \hfill(1.0)

\item[(T12.1c)] Power emitted by an evaporating black hole
  $= \sigma T_{\mathrm{bh}}^4 4\pi R_{\mathrm{SC}}^2$. \hfill(1.0)

  This gives
  \[
    \frac{dM_{\mathrm{bh}}}{dt}
    = -\frac{1}{c^2}\sigma T_{\mathrm{bh}}^4 4\pi R_{\mathrm{SC}}^2.
  \]
  \hfill(1.0)

  With $T_{\mathrm{bh}} = b/(16 R_{\mathrm{SC}})$,
  \[
    \frac{dM_{\mathrm{bh}}}{dt}
    = -\frac{\pi b^4 \sigma c^2}{16^4 G^2}\,\frac{1}{M_{\mathrm{bh}}^2}
  \]
  \hfill(1.0)

  This gives
  $\dfrac{dM_{\mathrm{bh}}}{M_{\mathrm{bh}}^{2}}
  = -\dfrac{\pi b^4 \sigma c^2}{16^4 G^2}\,dt$,
  % sic: source writes dM_bh/M_bh^2, though the integral that follows is of
  % M_bh^2 dM_bh
  and hence
  $\dfrac{M_{\mathrm{bh}}^3}{3}
  = -\dfrac{\pi b^4 \sigma c^2}{16^4 G^2}\,t + C$. \hfill(1.0)
  Since $M_{\mathrm{bh}} = M_0$ at $t = 0$, we get
  \[
    M_{\mathrm{bh}}^3 = M_0^3 - 3\frac{\pi b^4 \sigma c^2}{16^4 G^2} t
  \]
  \hfill(0.5)
  \[
    M_{\mathrm{bh}}(t)
    = \left(M_0^3 - 3\frac{\pi b^4 \sigma c^2}{16^4 G^2} t\right)^{1/3}
  \]
  \hfill(1.0)

  The sketch of $M_{\mathrm{bh}}(t)$ is a convex curve starting at $M_0$ on
  the $M_{\mathrm{bh}}(t)$-axis, with negative initial slope and large
  negative final slope, reaching zero at
  $t = \dfrac{16^4 G^2 M_0^3}{3\pi b^4 \sigma c^2}$. \hfill(2.5)
  % FIGURE NEEDED: answer sketch of M_bh(t) vs t, convex decreasing from
  % (0, M_0) to (16^4 G^2 M_0^3 / (3 pi b^4 sigma c^2), 0)
  % (2025 theory solutions, page 54 of 59)

\item[(T12.1d)] We know that
  $M(t) = \left(M_0^3 - 3\dfrac{\pi b^4 \sigma c^2}{16^4 G^2}
  t\right)^{1/3}$.

  Clearly, $M(t) = 0$ when
  $t = \dfrac{M_0^3}{3}\,\dfrac{16^4 G^2}{\pi b^4 \sigma c^2}$. Thus,
  \[
    \tau_{\mathrm{bh}} = \frac{M_0^3}{3}\,
      \frac{16^4 G^2}{\pi b^4 \sigma c^2}
  \]
  \hfill(1.0)

  With $A \equiv \dfrac{\pi b^4 \sigma c^2}{16^4 G^2}
  \simeq 3.868 \times 10^{15}\,\mathrm{kg^3\,s^{-1}}$, this gives \hfill(0.5)
  \[
    \tau_{\mathrm{bh}} \approx \frac{M_0^3}{3A}
    = 8.617 \times 10^{-17} \left(\frac{M_0}{\mathrm{kg}}\right)^3
    \mathrm{s}.
  \]
  \hfill(0.5)

  For $M_0 = 10\,\mathrm{M_\odot}$,
  \[
    \tau_{\mathrm{bh,\,10\odot}} \approx 6.8 \times 10^{77}\,\mathrm{s}
  \]
  \hfill(1.0)

\item[(T12.1e)] As the CMB will redshift away eventually, there will be no
  CMB radiation for the BH to absorb at large $t$. It will only decrease in
  mass and increase in temperature. So only the figures which show the
  temperature of the black hole increasing eventually ($t \to \infty$) are
  relevant. The derivatives at the origin (now) should distinguish the three
  cases.
  \begin{description}
  \item[(X)] If $T_{\mathrm{bh}}^{\mathrm{now}} >
    T_{\mathrm{cmb}}^{\mathrm{now}}$, then the BH is losing more energy now
    than it is gaining by CMB absorption, and its temperature should increase
    due to decreasing mass. $(dT/dt)_{\mathrm{now}} > 0$.
  \item[(Y)] If $T_{\mathrm{bh}}^{\mathrm{now}} =
    T_{\mathrm{cmb}}^{\mathrm{now}}$, the temperature will be in equilibrium
    for some time, i.e.\ $(dT/dt)_{\mathrm{now}} = 0$, before the CMB
    temperature redshifts to smaller values and $(dT/dt)_{\infty} > 0$.
  \item[(Z)] For $T_{\mathrm{bh}}^{\mathrm{now}} <
    T_{\mathrm{cmb}}^{\mathrm{now}}$, initially the BH will gain mass (and
    hence cool down) by absorbing from CMBR compared to what it loses by
    evaporation. So $(dT/dt)_{\mathrm{now}} < 0$. But eventually,
    $(dT/dt)_{\infty} > 0$.
  \end{description}
  Therefore: X $\to$ III, \quad Y $\to$ II, \quad Z $\to$ VI. \hfill(6.0)

\item[(T12.2a)] We want
  $\tau_{\mathrm{PBH}} = 14 \times 10^{9} \times 3.156 \times 10^{7}\,\mathrm{s}
  = 4.418 \times 10^{17}$\,s. \hfill(1.0)

  But we have found in (T12.1d) that
  $\tau_{\mathrm{bh}} = 8.617 \times 10^{-17}
  \left(\dfrac{M_0}{\mathrm{kg}}\right)^3$\,s. This gives
  \[
    M_{0,\mathrm{PBH}}^3
    = \left(\frac{\tau_{\mathrm{PBH}}}{8.617 \times 10^{-17}}\right)
      \mathrm{kg}^3
    \Rightarrow M_{0,\mathrm{PBH}} \simeq 1.7 \times 10^{11}\,\mathrm{kg}
  \]
  \hfill(1.0)

  Since $R_{\mathrm{SC}} \propto M_{\mathrm{bh}}$, and we know
  $R_{\mathrm{SC,\odot}} = 2952$\,m,
  \[
    R_{\mathrm{SC,\,PBH}} \simeq R_{\mathrm{SC,\odot}}
    \left(\frac{M_{0,\mathrm{PBH}}}{M_\odot}\right)
    \Rightarrow R_{\mathrm{SC,\,PBH}} = 2.6 \times 10^{-16}\,\mathrm{m}
  \]
  \hfill(1.0)

  Since $T_{\mathrm{bh}} \propto 1/M_{\mathrm{bh}}$ and we know
  $T_{\mathrm{bh,\,10\odot}}$ from (T12.1b),
  \[
    T_{\mathrm{PBH}} \simeq T_{\mathrm{bh,\,10\odot}}
    \left(\frac{10 M_\odot}{M_{0,\mathrm{PBH}}}\right)
    \Rightarrow T_{\mathrm{PBH}} = 7.1 \times 10^{11}\,\mathrm{K}
  \]
  \hfill(1.0)

\item[(T12.2b)] The size of horizon $r_{\mathrm{H}} = ct$. So the total mass
  contained within the horizon volume is
  \[
    M = \frac{1}{c^2}\rho \left(\frac{4\pi}{3}\right)
      \left(\frac{ct}{2}\right)^3.
  \]
  \hfill(2.0)

  With $a(t) \sim t^{1/2}$ and $a(t) \sim 1/T$, one gets
  $t \sim T^{-2}$. \hfill(2.0)

  With $\rho \sim T^4 \sim t^{-2}$, the mass of PBH formed at time $t$ would
  be
  \[
    M_{0,\mathrm{PBH}} \propto t.
  \]
  \hfill(1.0)

  The age of the Universe at which PBH of $10^{12}$\,kg forms is
  $t_{12} = 10^{-23}$\,s. Hence,
  \[
    t_{20} = t_{12} \times
    \left(\frac{10^{20}\,\mathrm{kg}}{10^{12}\,\mathrm{kg}}\right)
    \Rightarrow t_{20} = t_{12} \times 10^{8} = 10^{-15}\,\mathrm{s}
  \]
  \hfill(1.0)

\item[(T12.2c)] Lifetime of the PBH with
  $M_{0,\mathrm{PBH}} = 10^{10}$\,kg will be
  \[
    \tau_{\mathrm{PBH}} = 8.617 \times 10^{-17} \cdot (10^{10})^3\,\mathrm{s}
    = 8.617 \times 10^{13}\,\mathrm{s}.
  \]
  \hfill(0.5)

  As given in (T12.1b), $R_{\mathrm{SC,\odot}} = 2952$\,m. Using
  $R_{\mathrm{SC}} \propto M$, we get the peak wavelength of radiation
  emitted by this $10^{10}$\,kg PBH, at the time of emission, as
  \[
    \lambda_0 \simeq 16 R_{\mathrm{SC}}
    = 16 \cdot \frac{10^{10}\,\mathrm{kg}}{M_\odot} \cdot
      R_{\mathrm{SC,\odot}}
    = 2.4 \times 10^{-16}\,\mathrm{m}.
  \]
  \hfill(1.5)

  Redshifted wavelength received at earth at present time $t_0$ is obtained
  from
  \[
    \lambda_{\mathrm{earth}}
    = \lambda_0 \frac{a(t_0)}{a(\tau_{\mathrm{PBH}})}.
  \]
  \hfill(1.0)

  With $a(t) \sim t^{2/3}$, we get
  \[
    \lambda_{\mathrm{earth}}
    = \lambda_0 \left(\frac{t_{\mathrm{now}}}{\tau_{\mathrm{PBH}}}\right)^{2/3}
    = 2.4 \times 10^{-16}\,\mathrm{m}
      \left(\frac{4.418 \times 10^{17}\,\mathrm{s}}
        {8.617 \times 10^{13}\,\mathrm{s}}\right)^{2/3}
  \]
  \hfill(1.0)
  \[
    \lambda_{\mathrm{earth}} = 7.1 \times 10^{-14}\,\mathrm{m}
  \]
  \hfill(1.0)

\item[(T12.2d)] Using the expression for $T_{\mathrm{bh}}$ found in
  (T12.1b), the energy of a typical photon emitted at time $t$ for a PBH is
  \[
    E_{\mathrm{emitted}}(t) = k_{\mathrm{B}} T_{\mathrm{bh}}(t)
    = \frac{b c^2 k_{\mathrm{B}}}{32 G M_{\mathrm{bh}}(t)}.
  \]
  \hfill(1.0)

  From (T12.1c), we know
  \[
    M_0^3 = 3\frac{\pi b^4 \sigma c^2}{16^4 G^2} t + M_{\mathrm{bh}}(t)^3
    = 3At + \left(\frac{b c^2 k_{\mathrm{B}}}
      {32 G E_{\mathrm{emitted}}(t)}\right)^3.
  \]
  \hfill(1.0)

  The value of $E_{\mathrm{det}}$ can be different from
  $E_{\mathrm{emitted}}$ due to redshift. Hence, the smallest possible value
  of $E_{\mathrm{emitted}}(t)$ is
  $E_{\mathrm{det}} = 3.0 \times 10^{20}$\,eV. This gives the maximum value
  of the second term as
  \[
    \left(\frac{b c^2 k_{\mathrm{B}}}
      {32 G E_{\mathrm{emitted}}(t)}\right)^3_{\mathrm{max}}
    = 4.3 \times 10^{-5}\,\mathrm{kg}^3.
  \]
  \hfill(1.0)

  Given the value of $A = 3.868 \times 10^{15}\,\mathrm{kg^3\,s^{-1}}$, the
  second term can always be neglected for $t \gtrsim 10^{-19}$\,s. This is
  true in this problem. So $M_0 \approx (3At)^{1/3}$. \hfill(1.0)

  \textbf{The largest possible} $M_0$ would be obtained when $t$ is as large
  as possible, i.e., $t = t_{\mathrm{now}}$. This will correspond to a BH
  radiating in its last stages at present in the nearby region. \hfill(1.0)

  Then
  \[
    M_0^{\mathrm{max}} = (3At_{\mathrm{now}})^{1/3}
    = (5.13 \times 10^{33}\,\mathrm{kg^3})^{1/3}.
  \]
  \hfill(0.5)

  Thus,
  \[
    M_0^{\mathrm{max}} = 1.7 \times 10^{11}\,\mathrm{kg}.
  \]
  \hfill(0.5)

  Note that this was already found in (T12.2a), so the calculation above is
  not needed.

  \textbf{The smallest} $M_0$ will correspond to a BH that has evaporated as
  early as possible. Such a BH would have evaporated far away from Earth,
  and its radiation would be redshifted.

  The smallest time of evaporation would correspond to the largest redshift
  factor, and hence the largest $E_{\mathrm{emitted}}(t)$, which is
  $E_{\mathrm{emitted}}(t) = k_{\mathrm{B}} T_{\mathrm{Planck}}$.
  \hfill(1.0)

  The ratio of the scale factors at the time of the above PBH evaporation
  and now is
  \[
    \frac{a(t)}{a(t_{\mathrm{now}})}
    = \frac{E_{\mathrm{det}}}{k_{\mathrm{B}} T_{\mathrm{Planck}}}
    = \frac{3.0 \times 10^{20}\,\mathrm{eV}}{10^{19}\,\mathrm{GeV}}
    = 3.0 \times 10^{-8}
  \]
  With $a(t_{\mathrm{now}}) = 1.0$, we have
  $a(t) = 3.0 \times 10^{-8}$. \hfill(1.0)

  Using the plot of $a(t)$ vs $t$ from the figure, $t \simeq 10^{4}$\,s.
  Note that for this $t$, the second term in the original $M_0^3$ expression
  is still negligible. \hfill(1.0)

  Therefore,
  \[
    M_0^{\mathrm{min}} \approx (3At)^{1/3}
    \approx 4.9 \times 10^{6}\,\mathrm{kg}.
  \]
  \hfill(1.0)
\end{description}
